\begin{frame}{6.3 — O que é recursividade?}
\begin{block}{Definição}
Uma função recursiva resolve parte do problema chamando a si mesma com um problema menor.
\end{block}
\begin{enumerate}
\item \textbf{Caso-base}: resultado conhecido, sem nova chamada.
\item \textbf{Passo recursivo}: reduzir o problema.
\item \textbf{Progresso}: garantir que as chamadas alcancem o caso-base.
\end{enumerate}
\small Exemplo didático: $S(0)=0$ e $S(n)=n+S(n-1)$ para inteiros $n>0$. Na aula, limitaremos $n$ a 20.
\end{frame}
